%% kapa: exam=1 round=29 year=2018 period=2 form=A subject=appraisal_law subject_ko=감정평가관계법규 seq=22 number=22
%% source: Q-Net 감정평가사 기출문제 (cst00309) artlSeq=5202808 file=extracted/1008/5202808/2018년 제29회 감정평가사 1차시험 기출문제/2018년 제29회 감정평가사 1차시험 기출문제/2교시 A형.pdf
%% conversion: kordoc 4.24.51 -> tools/gen1.py (auto); review=proofread-sonnet; flags=none
\begin{question}{22}
국유재산법상 행정재산에 관한 설명으로 옳지 않은 것은?
\begin{choices}
  \item 행정재산의 사용허가를 철회하려는 경우에는 청문을 하여야 한다.
  \item 행정재산의 사용허가에 관하여는「국유재산법」에서 정한 것을 제외하고는 「민법」의 규정을 준용한다.
  \item 행정재산으로 할 목적으로 기부를 받은 재산에 대하여 중앙관서의 장이 기부자에게 사용허가하는 경우 그 사용료를 면제할 수 있다.
  \item 행정재산의 사용허가를 받은 자가 해당 재산의 보존을 게을리 한 경우 그 허가를 철회할 수 있다.
  \item 행정재산으로 사용하기로 결정한 날부터 5년이 지난 날까지 해당 재산이 행정재산으로 사용되지 아니한 경우 지체 없이 행정재산의 용도를 폐지하여야 한다.
\end{choices}
\end{question}
\answer{2}
